% Generated draft: requires mathematical and source review.
\documentclass[11pt]{article}
\usepackage[margin=25mm]{geometry}
\usepackage{fontspec}
\setmainfont{Latin Modern Roman}
\usepackage{amsmath,amssymb,mathtools,unicode-math}
\setmathfont{Latin Modern Math}
\usepackage{enumitem,needspace,xurl,hyperref}
\setlength{\parindent}{0pt}
\setlength{\parskip}{5pt}
\setlength{\emergencystretch}{3em}
\newcommand{\N}{\mathbb{N}}
\newcommand{\Z}{\mathbb{Z}}
\newcommand{\Q}{\mathbb{Q}}
\newcommand{\R}{\mathbb{R}}
\newcommand{\C}{\mathbb{C}}
\newcommand{\F}{\mathbb{F}}
\newcommand{\A}{\mathbb{A}}
\newcommand{\PP}{\mathbb{P}}
\newcommand{\E}{\mathbb{E}}
\newcommand{\eps}{\varepsilon}
\newcommand{\dd}{\,\mathrm{d}}
\begin{document}
\section*{UnsolvedMath ID: 20000803}
\textbf{Source problem code:} catalogue entry
\subsection{Definitions and mathematical statement}
Let $k$ be algebraically closed, $n,d\ge1$, and $R=k[x_1,\ldots,x_n]$. Let $H=\operatorname{Hilb}^d(\mathbb A_k^n)$ and let $H_{\mathrm{sm}}$ be the reduced closure of its locus of $d$ distinct points, the smoothable component. At a length-$d$ ideal $I$, the whole-Hilbert tangent space is $\operatorname{Hom}_R(I,R/I)$. The original asks to describe $T_{[I]}H_{\mathrm{sm}}$ for smoothable $I$, rather than assuming it equals this ambient tangent space at every point. The original does not specify the base field; the attached aid works in arbitrary characteristic over algebraically closed fields.

The aid treats the monogenic chart on which $1,x_1,\ldots,x_1^{d-1}$ is a basis of $R/I$. Its ideals have the form
\[
I=(f(x_1),x_2-p_2(x_1),\ldots,x_n-p_n(x_1)),
\]
where $f$ is monic of degree $d$ and $\deg p_i<d$. Their coefficient vectors form $\mathbb A_k^{nd}$. The chart lies in $H_{\mathrm{sm}}$, and at every point on it the aid gives
\[
T_{[I]}H_{\mathrm{sm}}=T_{[I]}H=\operatorname{Hom}_R(I,R/I)\simeq(R/I)^n,
\qquad\dim_k T_{[I]}H_{\mathrm{sm}}=nd.
\]
The displayed isomorphism uses the chosen generators or coefficient coordinates.

More precisely, for a coefficient vector $q$ and any tangent vector $v\in k^{nd}$, choose a vector $q_*$ with squarefree first polynomial and set
\[
\gamma(s)=q+s v+s^2(q_*-q-v).
\]
The aid states that this is a polynomial flat smoothing family in the chart, with $\gamma(0)=q$, $\gamma'(0)=v$, and generic fiber $d$ geometrically distinct points; indeed $\gamma(1)=q_*$ ensures the discriminant is not identically zero. Thus every chart tangent vector is realized by a smoothing with its exact prescribed first jet, in every characteristic. After a coordinate change the chart contains every curvilinear length-$d$ scheme, meaning each local component has embedding dimension at most one. The source supplies no corresponding description for all noncurvilinear points of the smoothable component.
\subsection{Short English statement}
Describe tangents to the smoothable Hilbert component. The supplied aid gives the full answer on the monogenic chart, including curvilinear schemes, and realizes every tangent there by an exact-first-jet smoothing. It does not identify all tangents at arbitrary smoothable points.
\subsection{Sources}
\begin{itemize}
\item Local UnsolvedMath registry, permanent ID $20000803$, source code catalogue entry; canonical \texttt{aim-arithmetic-geometry-notes.json}, \texttt{notes[48]}.
\item American Institute of Mathematics, workshop ``Components of Hilbert schemes,'' problem 2: \url{https://aimath.org/WWN/hilbertschemes/hilbertschemes.pdf}. The exact extracted question supplies no remarks or literature references. No PDF inspection is claimed here.
\item Attached research attempt 3: partial result, recorded run status open, rigor proved, novelty confidence low; registry status partially solved, classification PARTIAL-PROGRESS. It is a machine-generated, non-peer-reviewed aid requiring independent verification before citation. No authoritative reference for its prescribed-jet lemma is supplied.
\end{itemize}
\end{document}
